\chapter{Running each code: install, register, a first calculation}
\chaptermark{Running each code}
\label{ch:running-each-code}
\index{Codes!installing}
\index{Codes!first calculation}
\index{ORCA}
\index{Gaussian 16}
\index{NWChem}
\index{MOPAC}
\index{xTB}
\index{CREST}
\index{Quantum ESPRESSO}
\index{CP2K}
\index{GROMACS}
\index{AutoDock Vina}
\index{MLP (MACE)}
\index{ECCE-QM}
\eccexample{Examples > Getting Started}
\helppage{running-codes-on-linux}
\helppage{running-codes-on-macos}
\helppage{running-codes-on-windows}
\eccesource{help/src/running-codes-on-linux.md help/src/running-codes-on-macos.md help/src/running-codes-on-windows.md}
\checkedwith{9.5.3-dev}{4c8b4b851}

This appendix is the place to look up a code you have not yet used in ECCE. For each code it says how to install it and how to tell ECCE where it is, and then points to the chapter that takes you through a first calculation with it step by step. The general steps of registering a code are in Chapters~\ref{ch:machines} and \ref{ch:code-registration}. Where a code has no chapter of its own yet, its first calculation is the water optimisation of Chapter~\ref{ch:first-calculation}, with the changes listed in its section, so that every code can be tried on the same molecule and the results compared.

\section{ECCE-QM}
\label{sec:run-ecceqm}
\paragraph{Install and register.} Nothing to do: ECCE-QM comes with ECCE and is always offered on \texttt{localhost}, on Linux, macOS and Windows.

\paragraph{First calculation.} Chapter~\ref{ch:first-calculation}, which optimises water with ECCE-QM and shows its frequencies.

\section{xTB}
\label{sec:run-xtb}
\paragraph{Install and register.} \texttt{sudo apt install xtb} on Debian and Ubuntu. In \menu{Register Machines}, select \texttt{localhost}, then xTB on the \menu{Codes} tab, press \menu{Find} and \menu{Save}.

\paragraph{First calculation.} Chapter~\ref{ch:xtb}.

\section{CREST}
\label{sec:run-crest}
\paragraph{Install and register.} Download the Linux program from the CREST releases on GitHub (\texttt{crest-lab/crest}, version 3.0.2 or later), unpack it and copy the program \texttt{crest} to \texttt{/usr/local/bin}. CREST runs xTB itself, so install and register xTB first; then register CREST as for xTB.

\paragraph{First calculation.} Chapter~\ref{ch:crest}, which starts from an xTB optimisation.

\section{ORCA}
\label{sec:run-orca}
\paragraph{Install and register.} Download ORCA from FACCTS (\url{https://www.faccts.de/orca/}, free academic account) for your system and unpack it into a folder, for example \texttt{\textasciitilde/orca}. In \menu{Register Machines}, select ORCA on the \menu{Codes} tab and type the full path of the file \texttt{orca} (\texttt{orca.exe} on Windows) into \menu{Program}; never the bare name \texttt{orca}, which on Debian is the GNOME screen reader. For more than one core ORCA also needs MPI (Chapter~\ref{ch:code-registration}).

\paragraph{First calculation.} Chapter~\ref{ch:scans}, which optimises butadiene with ORCA and then scans its torsion. For the shortest first test, follow Chapter~\ref{ch:first-calculation} with ORCA instead of ECCE-QM: in step~1 make an ORCA calculation, and in the editor choose the theory \menu{RDFT}, then the functional B3LYP under \menu{Theory Details...}, the run type \menu{GeoVib} and the basis set def2-SVP.

\section{Gaussian 16}
\label{sec:run-g16}
\paragraph{Install and register.} A site with a licence installs Gaussian; ask for the path of \texttt{g16} and of its installation folder. Register the path in \menu{Program}, and add the line \texttt{g16root} followed by the folder above \texttt{g16} under \menu{Environment variables} (Chapter~\ref{ch:code-registration}).

\paragraph{First calculation.} Follow Chapter~\ref{ch:first-calculation} with these changes: in step~1 make a Gaussian~16 calculation; in the editor choose the theory \menu{RDFT}, the functional B3LYP under \menu{Theory Details...}, the run type \menu{GeoVib} and the basis set 6-31G*.

\section{NWChem}
\label{sec:run-nwchem}
\paragraph{Install and register.} \texttt{sudo apt install nwchem} on Debian; on a Mac from conda-forge with Miniforge. There is no Windows version: on Windows, register a Linux machine that has NWChem (Chapter~\ref{ch:hpc}). Register with \menu{Find}.

\paragraph{First calculation.} Follow Chapter~\ref{ch:first-calculation} with these changes: make an NWChem calculation, and in the editor choose the theory \menu{RHF}, the run type \menu{GeoVib} and the basis set 6-31G*.

\section{MOPAC}
\label{sec:run-mopac}
\paragraph{Install and register.} \texttt{sudo apt install mopac} on Debian. On macOS and Windows, run the installer from the MOPAC web site and install into a folder in your home folder, which needs no administrator rights; the program is \texttt{bin/mopac} (\texttt{bin\textbackslash mopac.exe}) in that folder. Register with \menu{Find} on Linux, otherwise with the full path.

\paragraph{First calculation.} Follow Chapter~\ref{ch:first-calculation} with these changes: make a MOPAC calculation, and in the editor choose the theory \menu{RPM7} and the run type \menu{GeoVib}. A semi-empirical method needs no basis set.

\section{Quantum ESPRESSO}
\label{sec:run-qe}
\paragraph{Install and register.} \texttt{sudo apt install quantum-espresso quantum-espresso-data-sssp} on Debian; the second package puts the SSSP pseudopotentials into \texttt{/usr/share/espresso/pseudo}. There is no Windows version. Register \texttt{pw.x} with \menu{Find}.

\paragraph{First calculation.} Chapter~\ref{ch:qe-settings}, which builds silicon from its space group and computes its energy.

\section{CP2K}
\label{sec:run-cp2k}
\paragraph{Install and register.} \texttt{sudo apt install cp2k} on Debian. Register with \menu{Find}.

\paragraph{First calculation.} Chapter~\ref{ch:cp2k}.

\section{GROMACS}
\label{sec:run-gromacs}
\paragraph{Install and register.} \texttt{sudo apt install gromacs} on Debian. There is no Windows version. Register \texttt{gmx} with \menu{Find}.

\paragraph{First calculation.} Chapter~\ref{ch:liquids}, which builds and runs a liquid.

\section{AutoDock Vina}
\label{sec:run-vina}
\paragraph{Install and register.} \texttt{sudo apt install autodock-vina} on Debian. Register \texttt{vina} with \menu{Find}.

\paragraph{First calculation.} Chapter~\ref{ch:docking}.

\section{MLP (MACE)}
\label{sec:run-mlp}
\paragraph{Install and register.} MACE needs a Python environment with PyTorch, which ECCE does not ship. On the \menu{Codes} tab, select MLP and press \menu{Set up...}, which builds the environment and downloads the default model (a download of a few gigabytes); then register the Python program of that environment, written out in full, as the MLP program.

\paragraph{First calculation.} Chapter~\ref{ch:mlp}.

